% Options for packages loaded elsewhere
\PassOptionsToPackage{unicode}{hyperref}
\PassOptionsToPackage{hyphens}{url}
\documentclass[
  11pt,
  a4paper,
]{scrartcl}
\usepackage{xcolor}
\usepackage[top=23mm,bottom=23mm,left=24mm,right=24mm]{geometry}
\usepackage{amsmath,amssymb}
\setcounter{secnumdepth}{-\maxdimen} % remove section numbering
\usepackage{iftex}
\ifPDFTeX
  \usepackage[T1]{fontenc}
  \usepackage[utf8]{inputenc}
  \usepackage{textcomp} % provide euro and other symbols
\else % if luatex or xetex
  \usepackage{unicode-math} % this also loads fontspec
  \defaultfontfeatures{Scale=MatchLowercase}
  \defaultfontfeatures[\rmfamily]{Ligatures=TeX,Scale=1}
\fi
\usepackage{lmodern}
\ifPDFTeX\else
  % xetex/luatex font selection
\fi
% Use upquote if available, for straight quotes in verbatim environments
\IfFileExists{upquote.sty}{\usepackage{upquote}}{}
\IfFileExists{microtype.sty}{% use microtype if available
  \usepackage[]{microtype}
  \UseMicrotypeSet[protrusion]{basicmath} % disable protrusion for tt fonts
}{}
\usepackage{setspace}
\makeatletter
\@ifundefined{KOMAClassName}{% if non-KOMA class
  \IfFileExists{parskip.sty}{%
    \usepackage{parskip}
  }{% else
    \setlength{\parindent}{0pt}
    \setlength{\parskip}{6pt plus 2pt minus 1pt}}
}{% if KOMA class
  \KOMAoptions{parskip=half}}
\makeatother
\usepackage{longtable,booktabs,array}
\newcounter{none} % for unnumbered tables
\usepackage{calc} % for calculating minipage widths
% Correct order of tables after \paragraph or \subparagraph
\usepackage{etoolbox}
\makeatletter
\patchcmd\longtable{\par}{\if@noskipsec\mbox{}\fi\par}{}{}
\makeatother
% Allow footnotes in longtable head/foot
\IfFileExists{footnotehyper.sty}{\usepackage{footnotehyper}}{\usepackage{footnote}}
\makesavenoteenv{longtable}
\setlength{\emergencystretch}{3em} % prevent overfull lines
\providecommand{\tightlist}{%
  \setlength{\itemsep}{0pt}\setlength{\parskip}{0pt}}

\usepackage{microtype}
\usepackage{xcolor}
\usepackage{needspace}
\definecolor{linkink}{HTML}{553653}
\AtBeginDocument{\hypersetup{colorlinks=true,linkcolor=linkink,urlcolor=linkink,citecolor=linkink}}
\setlength{\emergencystretch}{3em}
\setlength{\parindent}{0pt}
\setlength{\parskip}{0.42em}
\widowpenalty=10000
\clubpenalty=10000
\displaywidowpenalty=10000
\allowdisplaybreaks[1]
\usepackage{bookmark}
\IfFileExists{xurl.sty}{\usepackage{xurl}}{} % add URL line breaks if available
\urlstyle{same}
\hypersetup{
  pdftitle={Order transfer through six for Williamson low-storage Runge--Kutta methods},
  pdfauthor={Anonymous},
  hidelinks,
  pdfcreator={LaTeX via pandoc}}

\title{Order transfer through six for Williamson low-storage
Runge--Kutta methods}
\author{Anonymous}
\date{20 October 2026}

\begin{document}
\maketitle

\setstretch{1.04}
\section{Abstract}\label{abstract}

Bazavov conjectured that replacing the additive updates in a Williamson
low-storage Runge--Kutta method by exponential updates preserves its
classical order for orders three, four and five. We prove this statement
for every finite stage count and extend it through order six. The proof
rests on an all-degree identity which expresses certain chronological
products using only classical tableau data. If \(A,b\) are the
associated classical coefficients, then for arbitrary vectors \(u,v\)
and every \(m\geq1\), \[
 t_m(u,v,\ldots,v)=\frac1{m!}b^TR_v^{m-1}u,\qquad
 R_v=(b^Tv)I-\operatorname{diag}(Av)+\operatorname{diag}(v)A.
\] A division-free row induction proves the identity, including
degenerate coefficient lists. Its polarisation and the shuffle
identities determine every required ordered tensor through total weight
six. An integration homomorphism then shows that the classical order
conditions force the exact ordered-flow coefficients. A fourth-order
proof in conventional tableau notation gives an elementary entry point.
Exact symbolic reduction checks all 197 ordered-tree coefficients
through degree six. At weight seven the determining equations have three
residual directions. Three explicit additional scalar tests suffice for
seventh-order transfer, but their automatic satisfaction remains
unproved. No unrestricted-order claim is made.

\section{1. Problem and contribution}\label{problem-and-contribution}

Many differential equations evolve within a constrained set of matrices:
rotations remain orthogonal, for example, and quantum evolution remains
unitary. An exponential update respects such a group constraint when its
exponent lies in the corresponding Lie algebra. The difficulty is
preserving the accuracy of an ordinary Runge--Kutta method after making
this change. Matrix products do not commute, so merely reusing the same
coefficients is generally insufficient. This paper identifies a
structured family for which the additional order conditions follow
automatically through order six.

Consider a sufficiently smooth matrix equation \[
Y'(t)=F(Y(t))Y(t).
\] Here \(F(Y)\) lies in the Lie algebra of the matrix group under
consideration. A natural geometric method updates \(Y\) by
multiplication by exponentials. Its matrix products do not commute, so
classical Runge--Kutta order conditions generally require additional
conditions to imply the same order on the group (Owren, 2006).

Williamson's two-register formulation is \[
K_i=U_iK_{i-1}+h f(y_{i-1}),\qquad y_i=y_{i-1}+V_iK_i,\quad 1\leq i\leq s,
\] with \(K_0=0\). Bazavov's corresponding lift is \[
Z_i=U_iZ_{i-1}+hF(Y_{i-1}),\qquad
Y_i=\exp(V_iZ_i)Y_{i-1},\quad Z_0=0.
\tag{1}
\] The left-multiplication convention in (1) is fixed throughout. The
last exponential appears on the left of the product, while the
chronological order used in the pullback tensors runs from early rows to
late rows.

Bazavov (2022, Section 4.3) conjectured order transfer for the
higher-stage third-order and the fourth- and fifth-order methods.
Bazavov (2026) still described the general claim as conjectural. Lystad
(2026) proved the all-stage third-order case. We establish the
explicitly conjectured orders three, four and five for every finite
stage count, and extend the implication to order six. One tensor
identity handles every coefficient list, while a small finite uniqueness
argument accounts for the order boundary. The all-stage third-order
theorem is prior work and is not claimed as new.

The practical consequence is precise: a developer who already has a
Williamson tableau of classical order \(p\leq6\) obtains a
commutator-free Lie-group method of at least that order by applying (1).
The theorem does not compare error constants, stable step sizes, or
computational performance between particular methods.

More explicitly, two successive updates give \[
Y_2=\exp(V_2Z_2)\exp(V_1Z_1)Y_0.
\] The latest exponential is on the left. The two-register description
refers to the recurrence, not to all workspace needed to evaluate \(F\)
or apply an exponential. The theorem concerns exact exponential updates;
errors from an approximate exponential routine require their own
control.

The proof has four steps. The low-storage recurrence forces an
alternating row identity. That identity determines tensors whose
arguments after the first are repeated. Polarisation and shuffle
equations recover the remaining orderings through weight six. Finally, a
continuous integration model identifies these expressions with the
coefficients of the exact flow. Section 5 first shows the mechanism at
order four without tree notation.

For use in a solver, the input is a verified classical Williamson
coefficient list of order \(p\leq6\), and the output is the recurrence
(1) with consistency order at least \(p\). This removes the need to
solve additional noncommutative order constraints for that lift.
Constructing the classical list, selecting its error constants and
assessing its stability remain separate tasks. Yan (2017) supplies a
thirteen-stage example of nonlinear order five, although its order for
autonomous linear systems is eight; Bazavov (2025, Section 5.4)
identifies its Williamson form. Printed decimal coefficients are not an
exact algebraic certificate. We supply a rational fourth-order example
below, but no sixth-order coefficient construction.

\section{2. Coefficients and ordered
tensors}\label{coefficients-and-ordered-tensors}

Vectors \(e_i^T\) are the coordinate rows of \(\mathbb R^s\). Define
rows \[
w_0=0,\quad w_i=U_iw_{i-1}+e_i^T,
\qquad \gamma_i=V_iw_i,
\] \[
a_i=\sum_{j<i}\gamma_j,
\qquad b^T=\sum_{i=1}^s\gamma_i.
\tag{2}
\] The classical Butcher matrix \(A\) has row \(a_i\); it is strictly
lower triangular. Set \[
c=A\mathbf1,\quad C=b^T\mathbf1,\quad D=\operatorname{diag}(c),\quad R=CI-D+A.
\tag{3}
\] No order or consistency assumption is imposed in this section. In
particular, \(C\) need not equal one.

For vectors \(u_1,\ldots,u_m\), define the chronological tensor \[
t_m(u_1,\ldots,u_m)=
\sum_{1\leq i_1\leq\cdots\leq i_m\leq s}
\frac{\prod_{r=1}^m(\gamma_{i_r}u_r)}{\prod_{j=1}^s n_j!},
\tag{4}
\] where \(n_j\) counts the occurrences of \(j\) among the indices.
Equation (4) is Owren's tensor (24), with his chronological exponential
rows \(\beta_j\) equal to \(\gamma_j\). In particular, \[
t_2(u,v)=\sum_{i<j}(\gamma_i u)(\gamma_jv)
+\frac12\sum_i(\gamma_i u)(\gamma_i v).
\tag{5}
\] Let \(T_{2,r}\) denote the same tensor for the prefix
\(\gamma_1,\ldots,\gamma_{r-1}\). This prefix is precisely the
exponential representation of stage \(r\).

For two rows the definition becomes \[
t_2(u,v)=\tfrac12(\gamma_1u)(\gamma_1v)
+(\gamma_1u)(\gamma_2v)
+\tfrac12(\gamma_2u)(\gamma_2v).
\] The halves come from choosing two factors from one exponential
series. The mixed term chooses one factor from each row in early-to-late
order. Reversing chronology replaces that mixed term by
\((\gamma_2u)(\gamma_1v)\), which need not agree. This chronological
convention is the pullback convention corresponding to the
left-multiplication product displayed above.

\section{3. A row identity}\label{a-row-identity}

\textbf{Lemma 1.} For every pair of vectors \(u,v\), every row in (2)
satisfies \[
w_i\bigl((Av)\odot u-v\odot(Au)\bigr)
=(a_iv)(w_iu)-(w_iv)(a_iu),
\tag{6}
\] where \(\odot\) denotes componentwise multiplication.

\textbf{Proof.} The claim for \(i=1\) is immediate. Suppose it holds for
\(i-1\). The prefix changes by \(a_i=a_{i-1}+V_{i-1}w_{i-1}\). Replacing
\(a_{i-1}\) by \(a_i\) in the alternating expression on the right does
not change its value, because the two added products cancel. The
coordinate row \(e_i^T\) satisfies (6), since \((Au)_i=a_iu\) and
\((Av)_i=a_iv\). Multiply the previous-row identity by \(U_i\), add the
coordinate-row identity, and use \(w_i=U_iw_{i-1}+e_i^T\). No division
occurs. \(\square\)

For fixed \(v\), put \[
C_v=b^Tv,\qquad Q_i=\gamma_iv,\qquad
 d_i=C_v-a_iv,\qquad d_{s+1}=0,
\] \[
R_v=C_vI-\operatorname{diag}(Av)+\operatorname{diag}(v)A.
\tag{7}
\] Then \(d_i-d_{i+1}=Q_i\). Multiplying (6) by \(V_i\) gives the row
identity \[
\gamma_iR_v=d_i\gamma_i+Q_i\sum_{j<i}\gamma_j.
\tag{8}
\] Equivalently, \(\Gamma R_v=K\Gamma\), where \(\Gamma\) has rows
\(\gamma_i\) and \[
K_{ij}=\begin{cases}Q_i&j<i,\\d_i&j=i,\\0&j>i.\end{cases}
\tag{9}
\] Neither matrix needs to be invertible. Taking \(v=\mathbf1\) recovers
\(R_v=R\) in (3).

\section{4. An all-degree contraction
identity}\label{an-all-degree-contraction-identity}

For \(k\geq0\), write \[
H_k(x,y)=\sum_{r=0}^k x^{k-r}y^r.
\] Thus \((x-y)H_k(x,y)=x^{k+1}-y^{k+1}\), including when \(x=y\).

\textbf{Lemma 2.} For every integer \(m\geq1\), \[
(\mathbf1^TK^{m-1})_i=H_{m-1}(d_i,d_{i+1}).
\tag{10}
\]

\textbf{Proof.} The case \(m=1\) is immediate. Multiplication by \(K\)
gives at index \(i\) \[
d_iH_{m-1}(d_i,d_{i+1})+
\sum_{j>i}Q_jH_{m-1}(d_j,d_{j+1}).
\] The sum telescopes to \(d_{i+1}^{m}\), because \(d_{s+1}=0\). Adding
this term to the first expression gives \(H_m(d_i,d_{i+1})\).
\(\square\)

\textbf{Theorem 3.} For every pair of real vectors \(u,v\), every finite
real Williamson coefficient list, and every integer \(m\geq1\), \[
\boxed{t_m(u,v,\ldots,v)=\frac1{m!}b^TR_v^{m-1}u.}
\tag{11}
\]

\textbf{Proof.} In (4), group the terms by the first index \(i\) and its
number of repetitions \(r\). The remaining arguments are \(v\), so
summing the later rows gives \[
t_m(u,v,\ldots,v)
=\sum_i(\gamma_iu)\sum_{r=1}^m
\frac{Q_i^{r-1}d_{i+1}^{m-r}}{r!(m-r)!}.
\tag{12}
\] The binomial theorem, interpreted as a polynomial identity when
\(Q_i=0\), identifies the inner sum as \(H_{m-1}(d_i,d_{i+1})/m!\).
Equations (8) and (10), together with \(b^T=\mathbf1^T\Gamma\), yield \[
b^TR_v^{m-1}=\mathbf1^TK^{m-1}\Gamma
=\sum_iH_{m-1}(d_i,d_{i+1})\gamma_i.
\] Substitute into (12). \(\square\)

The bilinear case is particularly simple: \[
2t_2(u,v)=(b^Tu)(b^Tv)
+b^T\bigl((Au)\odot v-(Av)\odot u\bigr).
\tag{13}
\] For \(m=2,u=c,v=\mathbf1\), the formula recovers Lystad's third-order
moment identity. Bazavov's c-reflection work contains the related
operator \(I-D+A\). These are antecedents of the present argument; we do
not claim novelty for that operator or for the earlier third-order
result.

\section{5. Fourth-order transfer}\label{fourth-order-transfer}

The eight classical conditions through order four are \[
\begin{array}{llll}
b^T\mathbf1=1,&b^Tc=\tfrac12,&b^Tc^2=\tfrac13,&b^TAc=\tfrac16,\\
b^Tc^3=\tfrac14,&b^T\operatorname{diag}(c)Ac=\tfrac18,
&b^TAc^2=\tfrac1{12},&b^TA^2c=\tfrac1{24}.
\end{array}
\tag{C4}
\] Powers of vectors in (C4) are componentwise. Owren (2006, equations
26a, 27a, 28a, 29a and 31a) proves that, after imposing (C4), the
remaining conditions through order four are \[
t_2(c,\mathbf1)=\tfrac16,\quad
 t_3(c,\mathbf1,\mathbf1)=\tfrac1{24},\quad
 t_2(c^2,\mathbf1)=\tfrac1{12},\quad
 t_2(Ac,\mathbf1)=\tfrac1{24},\quad
 \sum_rb_rT_{2,r}(c,\mathbf1)=\tfrac1{24}.
\tag{L4}
\] The first is the order-three condition. Its other orientation follows
from \(t_2(c,\mathbf1)+t_2(\mathbf1,c)=(b^Tc)(b^T\mathbf1)\).

\textbf{Theorem 4.} Every finite-stage Williamson method satisfying (C4)
has a lift (1) of consistency order at least four.

\textbf{Proof.} Since \(C=1\), apply (11) with \(m=2\) and
\(u=c,c^2,Ac\). The resulting values are respectively \[
\frac12(\tfrac12-\tfrac13+\tfrac16)=\tfrac16,
\quad\frac12(\tfrac13-\tfrac14+\tfrac1{12})=\tfrac1{12},
\quad\frac12(\tfrac16-\tfrac18+\tfrac1{24})=\tfrac1{24}.
\] For the third-order tensor, expand \((I-D+A)^2\), keeping the matrix
order: \[
6t_3(c,\mathbf1,\mathbf1)
=b^T\bigl(c-2c^2+2Ac+c^3-D Ac-Ac^2+A^2c\bigr).
\] Equation (C4) gives \[
6t_3(c,\mathbf1,\mathbf1)
=\tfrac12-\tfrac23+\tfrac13+\tfrac14-\tfrac18-\tfrac1{12}+\tfrac1{24}
=\tfrac14.
\] Finally apply the \(m=2\) formula separately at every stage prefix.
Its total mass is \(c_r\), so \[
2T_{2,r}(c,\mathbf1)=c_r(Ac)_r-(Ac^2)_r+(A^2c)_r.
\] Weighting with \(b_r\) gives \[
\sum_rb_rT_{2,r}(c,\mathbf1)
=\frac12(\tfrac18-\tfrac1{12}+\tfrac1{24})=\tfrac1{24}.
\] Thus every condition in (L4) holds. Owren's complete order theory
implies consistency order four. \(\square\)

Under the standard local regularity assumptions for one-step
convergence, the local defect is \(O(h^5)\) and the global error on a
fixed finite integration interval is \(O(h^4)\). For nonautonomous
equations, adjoining time gives the usual autonomous formulation; the
time component of the lifted constant frame reproduces the classical
stage nodes.

\section{6. Transfer through order
six}\label{transfer-through-order-six}

The fourth-order argument illustrates the mechanism without trees. To
handle orders five and six systematically, we combine Theorem 3 with the
shuffle relations and the elementary-weight form of the classical order
conditions.

\subsection{6.1 Classical and ordered
trees}\label{classical-and-ordered-trees}

Let \(\tau\) be an ordinary, unordered rooted tree. Its number of
vertices is \(|\tau|\), and its tree factorial is \[
\gamma(\bullet)=1,\qquad
\gamma([\tau_1,\ldots,\tau_k])
=|[\tau_1,\ldots,\tau_k]|\prod_{j=1}^k\gamma(\tau_j).
\] Its stage elementary vector is \[
\Phi_\bullet=\mathbf1,\qquad
\Phi_{[\tau_1,\ldots,\tau_k]}
=\bigodot_{j=1}^k A\Phi_{\tau_j},
\] and its output moment is \(M_\tau=b^T\Phi_\tau\). The classical
order-p conditions are \[
M_\tau=\frac1{\gamma(\tau)}\qquad (|\tau|\leq p).
\tag{C}
\] Every moment has weight \(|\tau|\); weights add when moments
multiply.

An ordered tree \(t=(t_1,\ldots,t_k)\) retains the order of its
children. Write \(g(t)\) for its numerical operator-series coefficient;
its degree is \(|t|-1\), so the one-vertex ordered tree represents the
identity operator. Owren's expansion gives \[
g(t)=t_k\bigl(g_{\mathrm{st}}(t_1),\ldots,g_{\mathrm{st}}(t_k)\bigr),
\tag{O}
\] where \(g_{\mathrm{st}}(t_j)\) is the vector of coefficients at the
evaluation-stage prefixes. The exact coefficients satisfy \[
e(\bullet)=1,\qquad
e((t_1,\ldots,t_k))
=\frac{\prod_{j=1}^k e(t_j)}{
\prod_{j=1}^k\left(\sum_{\ell=1}^j|t_\ell|\right)}.
\tag{E}
\] Equation (E) is the iterated-integral form of Owren's exact-flow
recursion (9). Our output \(g(t)\) is his \(g_{s+1,J}(t)\), and \(e(t)\)
is his \(\alpha(t)/(|t|-1)!\). Owren's Theorem 2.5 and Corollary 2.6
imply that matching \(g(t)=e(t)\) for \(|t|\leq p+1\) suffices for
consistency order \(p\) throughout the stated problem class. Necessity
is a universal order statement, not a claim about each individual
equation, where special cancellations can occur.

The two tree conventions serve different purposes:

\Needspace{13\baselineskip}

{\def\LTcaptype{none} % do not increment counter
\begin{longtable}[]{@{}
  >{\raggedright\arraybackslash}p{(\linewidth - 4\tabcolsep) * \real{0.3000}}
  >{\raggedleft\arraybackslash}p{(\linewidth - 4\tabcolsep) * \real{0.4000}}
  >{\raggedright\arraybackslash}p{(\linewidth - 4\tabcolsep) * \real{0.3000}}@{}}
\toprule\noalign{}
\begin{minipage}[b]{\linewidth}\raggedright
Object
\end{minipage} & \begin{minipage}[b]{\linewidth}\raggedleft
Bookkeeping degree
\end{minipage} & \begin{minipage}[b]{\linewidth}\raggedright
Meaning
\end{minipage} \\
\midrule\noalign{}
\endhead
\bottomrule\noalign{}
\endlastfoot
Ordinary tree \(\tau\) & \(|\tau|\) & Weight of the classical moment
\(M_\tau\) \\
Ordered operator tree \(t\) & \(|t|-1\) & Power of the step size in its
output coefficient \\
Prefix conversion of an output polynomial of weight \(d\) & \(d+1\) &
Weight of its stage elementary vectors \\
\end{longtable}
}

The singleton ordered tree represents the identity operator, so it has
degree zero. This accounts for the one-vertex shift; it is not an
additional order condition.

\subsection{6.2 Which tensors are
determined?}\label{which-tensors-are-determined}

The tensors (4) satisfy the shuffle identity \[
t_r(u_1,\ldots,u_r)t_q(v_1,\ldots,v_q)
=\sum_{\sigma\in\mathrm{Sh}(r,q)}t_{r+q}(\sigma(u_1,\ldots,u_r;v_1,\ldots,v_q)).
\tag{S}
\] Repeated words are counted with multiplicity. This follows either by
partitioning the ordered integration simplex or by multiplying the
defining products of exponential series.

Polarising Theorem 3 in its repeated vector supplies another relation.
For a fixed first argument \(u\), it determines the sum of all distinct
permutations of any prescribed tail multiset. The right side is \(1/m!\)
times the sum of the corresponding operator products in the linear maps
\(R_v\), contracted with \(b^T\). We denote this known equation by
\(P(u)\).

For example, the tail multiset \(\{a,a,b\}\) gives \[
t_4(u,a,a,b)+t_4(u,a,b,a)+t_4(u,b,a,a)
=\frac1{24}b^T(R_a^2R_b+R_aR_bR_a+R_bR_a^2)u.
\] Extracting the coefficient of \(x^2y\) after substituting \(v=xa+yb\)
gives exactly the three distinct words displayed. Equivalently, mixed
differentiation introduces a common factor \(2!\), which cancels on the
two sides. In contrast, repeated interleavings in a shuffle equation
retain their multiplicity.

\textbf{Lemma 5 (finite tensor determination).} For arguments chosen
among the ordinary elementary vectors \(\Phi_\tau\), with each argument
assigned weight \(|\tau|\), the polarised identities \(P(u)\) and the
shuffle identities determine every tensor of total weight at most six
from classical moment polynomials.

\textbf{Proof.} Use induction on tensor length. A tensor of length one
is a classical moment. A tensor with every argument equal follows from
\(t_m(v,\ldots,v)=(b^Tv)^m/m!\). For a prescribed multiset of arguments,
regard the tensors for its distinct permutations as unknowns. In each
shuffle equation, the factors on the right have smaller length and are
already known. The right side of each \(P(u)\) is a classical moment
polynomial because \(R_v\) uses only \(A\), \(b^T\), and componentwise
multiplication.

There is only one ordinary tree of weight one and only one of weight
two. Consequently the nonconstant multiplicity patterns that can occur
at total weight at most six are exactly those in the following table.
Letters denote distinct arguments, independently of their order in the
alphabet.

\Needspace{22\baselineskip}

{\def\LTcaptype{none} % do not increment counter
\begin{longtable}[]{@{}
  >{\raggedright\arraybackslash}p{(\linewidth - 6\tabcolsep) * \real{0.2143}}
  >{\raggedleft\arraybackslash}p{(\linewidth - 6\tabcolsep) * \real{0.2857}}
  >{\raggedright\arraybackslash}p{(\linewidth - 6\tabcolsep) * \real{0.2143}}
  >{\raggedleft\arraybackslash}p{(\linewidth - 6\tabcolsep) * \real{0.2857}}@{}}
\toprule\noalign{}
\begin{minipage}[b]{\linewidth}\raggedright
Pattern
\end{minipage} & \begin{minipage}[b]{\linewidth}\raggedleft
Unknown permutations
\end{minipage} & \begin{minipage}[b]{\linewidth}\raggedright
A square set of known equations
\end{minipage} & \begin{minipage}[b]{\linewidth}\raggedleft
Determinant
\end{minipage} \\
\midrule\noalign{}
\endhead
\bottomrule\noalign{}
\endlastfoot
ab & 2 & P(a), P(b) & 1 \\
aab & 3 & P(a), P(b), S(a,ab) & 1 \\
abc & 6 & P(a), P(b), P(c), S(a,bc), S(ab,c), S(ac,b) & −1 \\
aaab & 4 & P(a), P(b), S(a,aab), S(aa,ab) & 1 \\
aabb & 6 & P(a), P(b), S(a,abb), S(aab,b), S(a,bab), S(ab,ab) & 2 \\
aaaab & 5 & P(a), P(b), S(a,aaab), S(aa,aab), S(aaa,ab) & 1 \\
\end{longtable}
}

Here S(x,y) means the shuffle equation for the two indicated words. The
integer coefficient matrices, with columns in lexicographic word order,
are printed in Appendix A. Each is invertible over the rationals. This
proves the induction. \(\square\)

The assertion that the right side is a classical moment polynomial can
be made explicit. If \(\tau=[\tau_1,\ldots,\tau_r]\) and
\(\sigma=[\sigma_1,\ldots,\sigma_q]\), then \[
R_{\Phi_\sigma}\Phi_\tau
=M_\sigma\Phi_\tau
-\Phi_{[\tau_1,\ldots,\tau_r,\sigma]}
+\Phi_{[\sigma_1,\ldots,\sigma_q,\tau]}.
\tag{G}
\] This formula proves closure. Counting the weight of a scalar moment
plus the weight of its stage tree, each term has weight
\(|\tau|+|\sigma|\). Applying \(R_{\Phi_\sigma}\) therefore raises stage
weight by \(|\sigma|\). Repeated applications of (G), followed by
contraction with \(b^T\), expand every polarised right side into a
rational polynomial in ordinary moments of the same total weight.

A polynomial in prefix moments is converted into a stage vector by
replacing each monomial \(M_{\tau_1}\cdots M_{\tau_k}\) with
\(\Phi_{[\tau_1,\ldots,\tau_k]}\). Indeed, its i-th component is
\(\prod_j(a_i\Phi_{\tau_j})\), exactly the same monomial at stage i. The
empty monomial becomes \(\Phi_\bullet=\mathbf1\). Thus (O), Theorem 3
and Lemma 5 express each ordered coefficient of degree at most six as a
rational polynomial in the classical moments. The polynomial for
\(g(t)\) is homogeneous of weight \(|t|-1\); converting it to a stage
vector gives arguments of weight \(|t|\).

For a concrete example, the ordered tree with two singleton children has
output coefficient \(t_2(\mathbf1,\mathbf1)=\tfrac12M_\bullet^2\). At
evaluation stage \(i\), this becomes \(\tfrac12c_i^2\), or
\(\tfrac12\Phi_{[\bullet,\bullet]}\): output weight two becomes stage
weight three. Truncating the recurrence after stage \(i-1\) leaves every
earlier component of an elementary vector unchanged because \(A\) is
strictly lower triangular. Therefore the prefix substitution is valid at
every stage, including degenerate coefficient lists.

\subsection{6.3 Why the resulting coefficients are the exact
ones}\label{why-the-resulting-coefficients-are-the-exact-ones}

We now evaluate those moment polynomials under (C). This evaluation has
a useful continuous interpretation. Replace the classical operations by
\[
(Af)(x)=\int_0^x f(y)\,dy,\qquad
b^Tf=\int_0^1f(x)\,dx,
\] and interpret componentwise multiplication as ordinary multiplication
of functions. The elementary vector then becomes the polynomial \[
\Phi_\tau(x)=\frac{|\tau|}{\gamma(\tau)}x^{|\tau|-1},
\] so \(b^T\Phi_\tau=1/\gamma(\tau)\). This is precisely the
substitution in (C), and it preserves every operation used in the
reduction.

The continuous ordered tensors are iterated integrals. They satisfy (S)
and the same repeated-tail formula. To verify the latter directly, write
\[
\rho(x)=\int_x^1v(y)\,dy,
\qquad (R_vu)(x)=\rho(x)u(x)+v(x)\int_0^xu(y)\,dy.
\] Changing integration order gives, for every nonnegative integer j, \[
\int_0^1\rho(x)^j(R_vu)(x)\,dx
=\frac{j+2}{j+1}\int_0^1\rho(x)^{j+1}u(x)\,dx.
\] Iterating yields \[
\frac1{m!}\int_0^1R_v^{m-1}u(x)\,dx
=\frac1{(m-1)!}\int_0^1u(x)\rho(x)^{m-1}\,dx,
\] which is the iterated integral with first argument u and m−1
subsequent arguments v. No positivity or monotonicity assumption is
needed.

\textbf{Theorem 6 (order transfer through six).} Let \(1\leq p\leq6\).
For every finite stage count and every real Williamson coefficient list
satisfying the classical order conditions (C), the lift (1) has
consistency order at least p.~In particular, it resolves the order-3,
order-4 and order-5 claim explicitly posed in Bazavov's Section 4.3 and
extends the implication to order six.

\textbf{Proof.} Induct over ordered-tree size in (O). The stage
coefficient for a child is the prefix version of the classical moment
polynomial already constructed for that child. A prefix moment evaluates
continuously to \(x^{|\tau|}/\gamma(\tau)\). Homogeneity therefore makes
the child stage coefficient \(e(t_j)x^{|t_j|-1}\). Lemma 5 determines
the root tensor uniquely using repeated-tail and shuffle equations.
Continuous iterated integrals satisfy those same equations. Hence
evaluation under (C) gives the corresponding continuous iterated
integral, whose value is exactly (E). Every needed moment has weight at
most \(|t|-1\leq p\), so all substitutions are warranted by the assumed
classical order conditions. Thus \(g(t)=e(t)\) for every
\(|t|\leq p+1\). Owren's completeness theorem gives consistency order
p.~\(\square\)

For the analytic convergence consequence, suppose the exact solution and
all sufficiently small-step numerical trajectories lie in a common
neighbourhood on a fixed interval \([0,T]\), with bounded derivatives
through the orders needed for the expansion and local Lipschitz bounds
there. The smooth one-step map then has a Lipschitz bound \(1+Lh\), and
the local defect is uniformly \(O(h^{p+1})\). Iterating
\(e_{n+1}\leq(1+Lh)e_n+Kh^{p+1}\), with exact initial data and
\(nh\leq T\), gives global error \(O(h^p)\). These analytic assumptions
are separate from the algebraic order conditions. The same estimate
permits an inexact step map whose uniform difference from the exact
lifted step is \(O(h^{p+1})\); it does not cover arbitrary exponential
approximations. In the nonautonomous case, adjoining an additive time
coordinate with constant vector field reproduces the classical stage
times on the extended state space.

\subsection{6.4 A precise boundary at weight
seven}\label{a-precise-boundary-at-weight-seven}

The tensor multiset \((a,a,a,b,b)\) has ten permutations. The complete
collection of its polarised repeated-tail and proper shuffle equations
has rank nine. With a the weight-one argument \(\mathbf1\) and b the
weight-two argument c, its total weight is seven. Appendix A gives a
nonzero null vector and a nonzero minor. Our identities therefore leave
one tensor direction undetermined at the next order. The script reports
this rank deficiency and stops. It is a boundary of this proof method,
not a counterexample to seventh-order transfer.

\subsection{6.5 Three explicit tests at order
seven}\label{three-explicit-tests-at-order-seven}

The missing direction just exhibited is not the only new family at
weight seven. A full enumeration gives two further deficient families,
involving the two ordinary vectors of weight three. This leads to a
conditional extension which makes the next-order problem finite and
explicit.

\textbf{Theorem 7 (a sufficient order-seven criterion).} Suppose a
finite real Williamson coefficient list satisfies all classical
order-seven conditions. If \[
t_5(\mathbf1,\mathbf1,c,\mathbf1,c)=\frac1{280},\qquad
t_4(\mathbf1,c,\mathbf1,c^2)=\frac1{84},\qquad
t_4(\mathbf1,c,\mathbf1,Ac)=\frac1{168},
\tag{C7+}
\] then its lift (1) has consistency order at least seven.

\textbf{Proof.} Every child of a degree-seven ordered coefficient has
degree at most six. Its prefix coefficient is therefore a linear
combination of ordinary elementary vectors by the preceding reduction.
The root tensors have total weight seven. Every factor of a proper
shuffle product has strictly smaller total weight and is already
determined. The polarised repeated-tail right sides are classical moment
polynomials of weight seven. After imposing the classical order-seven
conditions, the difference between the actual tensor values and their
continuous counterparts consequently lies in the nullspace of the
corresponding determination matrix.

There are \(1,1,2,4,9,20,48\) ordinary rooted trees of weights one
through seven. Enumerating multisets of these arguments with total
weight seven gives the following complete list; families distinguish the
actual ordinary-tree arguments, whereas patterns retain only
multiplicities.

\Needspace{30\baselineskip}

{\def\LTcaptype{none} % do not increment counter
\begin{longtable}[]{@{}lrr@{}}
\toprule\noalign{}
Pattern & Number of families & Rank / number of words \\
\midrule\noalign{}
\endhead
\bottomrule\noalign{}
\endlastfoot
a & 48 & 1/1 \\
ab & 37 & 2/2 \\
aab & 13 & 3/3 \\
abc & 5 & 6/6 \\
aaab & 5 & 4/4 \\
aaaab & 2 & 5/5 \\
aaaaab & 1 & 6/6 \\
aaaaaaa & 1 & 1/1 \\
aaabb & 1 & 9/10 \\
aabc & 2 & 11/12 \\
\end{longtable}
}

These counts can be reproduced by listing the partitions of seven and
assigning ordinary trees to each part. The matrices use exactly the
first-letter polarisation rows and all proper-shuffle rows defined in
Lemma 5. The full-rank cases through length five are already covered
there. For aaaaab, the rows
\(P(a),P(b),S(a,aaaab),S(aa,aaab),S(aaa,aab),S(aaaa,ab)\) give an
invertible six-by-six system. Constant-word tensors are determined
directly.

The aaabb family is \((\mathbf1,\mathbf1,\mathbf1,c,c)\). Its null
vector is printed in Appendix A and has nonzero aabab coordinate. The
first condition in (C7+) therefore removes its one free direction. The
aabc pattern occurs twice, as \((\mathbf1,\mathbf1,c,c^2)\) and
\((\mathbf1,\mathbf1,c,Ac)\). In lexicographic word order \[
(aabc,aacb,abac,abca,acab,acba,baac,baca,bcaa,caab,caba,cbaa)
\] its nullspace is spanned by \[
(0,0,-1,1,1,-1,1,-1,0,-1,1,0)^T.
\] Exact elimination of the polarisation and shuffle matrix gives rank
eleven. The abac coordinate is nonzero, so the remaining two conditions
remove both free directions. The small matrices and this exhaustive
family count are independently checkable by
\texttt{code/order7\_criterion.py}.

The constants follow directly from the continuous model. The iterated
integral of \(x_i^{w_i-1}\) over \(0<x_1<\cdots<x_m<1\) is \[
\left[\prod_{j=1}^m(w_1+\cdots+w_j)\right]^{-1}.
\] Weights \((1,1,2,1,2)\) give \(1/280\), and weights \((1,2,1,3)\)
give \(1/84\). Since \(Ac=x^2/2\) in that model, the third value is
\(1/168\). Every root tensor now agrees with its exact continuous value,
and the ordered-tree sufficiency argument proves the result. \(\square\)

At the level of matching the complete ordered series, the three tests
are also necessary: the required child-stage coefficients are
\(\mathbf1,c,c^2/2,Ac\), so the three tensors occur as root
coefficients, with a factor \(1/2\) in the second case. This is a
coefficient-level observation, not a necessity assertion for every
special differential equation. We have not proved that (C7+) follows
from the classical order-seven conditions, that its three conditions are
independent on genuine Williamson lists, or that a seventh-order list
exists. The unconditional transfer theorem remains bounded by six.

\section{7. Exact checks and scope}\label{exact-checks-and-scope}

The exact symbolic certificate records all 197 ordered-tree coefficients
through degree six, including the degree-zero identity coefficient, as
polynomials in ordinary moments. Its generator uses rational algebra,
the polarised form of Theorem 3 and shuffle equations; it uses no
numerical fitting. The independent checker imports only the Python
standard library and does not import the generator. It reads the emitted
formulas and verifies all 87 polarisation equations, 212 shuffle
equations, 47 nontrivial family ranks, 197 root-recursion substitutions
and 197 exact-flow values. It also independently enumerates the six
allowed multiplicity patterns and checks the matrices in Appendix A.
Thus the certificate can be checked without trusting the solver that
produced it.

A separate verifier reconstructs the tableau directly from the
low-storage recurrence. It performs 1,360 exact rational checks of the
row and tensor identities on 20 coefficient lists, including degenerate
lists and lists with negative coefficients, with tensor degree up to
eight. It independently expands the stage exponentials on four arbitrary
rational coefficient lists and compares all 197 emitted coefficient
polynomials, giving 788 further exact comparisons. It also checks the
eight classical and five ordered conditions for the rational five-stage
fourth-order method in Bazavov (2025, Table 7): \[
U=(0,-1,-1,-11,\tfrac1{10}),\qquad
V=(\tfrac12,\tfrac23,-\tfrac12,-\tfrac1{10},\tfrac16).
\tag{M4}
\] The direct ordered-tree expansion verifies every coefficient through
degree four for (M4), with the first defects at degree five. A nonlinear
noncommuting matrix check uses \[
F(Y)=X+Y_{12}W,\qquad
X=\begin{pmatrix}0&1\\-1&0\end{pmatrix},\qquad
W=\begin{pmatrix}1&1\\0&-1\end{pmatrix},\qquad Y(0)=I.
\] Direct truncated power-series arithmetic gives the first local-error
term for (M4) at degree five. Reversing the multiplication convention
gives a nonzero defect already at degree three, providing a negative
control. These computations check transcription, orientation and
concrete consequences; the all-stage theorem rests on the proofs in
Sections 3--6.

The theorem imposes no positivity, monotone-node, distinct-node, or
nonzero-coefficient requirement. It therefore covers both the regular
coefficient parameterisations in the source literature and their
direct-recurrence degeneracies. It does not assert that all classical
Runge--Kutta methods admit Williamson storage. The theorem is an
implication for every coefficient list satisfying the stated classical
order conditions. It makes no existence claim for a prescribed order and
stage count. The general repeated-tail tensor identity holds beyond
order six, but those identities and shuffle equations alone cease to
determine every required tensor at weight seven.

\section{8. Prior-art boundary}\label{prior-art-boundary}

The search, completed on 19 October 2026, inspected the originating
conjecture, Owren's complete order conditions, Bazavov's 2025
c-reflection preprint and 2026 coefficient paper, and Lystad's August
2026 proof. Exact formulas, the operator \(I-D+A\), low-storage aliases,
Lie--Butcher terminology, and recent resolution searches were checked.
The all-stage transfer theorem at orders four, five and six and the
all-degree repeated-tail contraction identity were not found in that
bounded corpus. This is evidence about the inspected corpus, not a proof
of historical priority. The detailed queries and
contribution-by-contribution decisions appear in
\texttt{novelty\_report.md}.

\section{9. Conclusion and
availability}\label{conclusion-and-availability}

The Williamson recurrence turns a restricted family of noncommutative
tensor products into classical moment data at every degree. Through
weight six, the remaining products are uniquely determined by
polarisation and shuffles, proving automatic transfer of the stated
classical orders. At weight seven the reduction leaves three scalar
tests. Establishing whether those tests are automatic, or finding a
genuine violating coefficient list, remains open.

The principal result is the written theorem, not a numerical trend. The
exact certificate checks transcription and finite algebra; the rational
and noncommuting examples check conventions and selected consequences.
None establishes external peer review, formal proof-assistant
verification, historical priority, or a performance advantage.

\textbf{Data and code availability.} The accompanying versioned
source-and-evidence package contains this manuscript, the order-six
certificate, the order-seven finite criterion checker, the independently
implemented checks and \texttt{AI\_INDEX.md}, which identifies the
replay entry point and the scope of every claim. The source repository
is \url{https://github.com/ipitchford/williamson-order-transfer} and the
versioned archive is \url{https://doi.org/10.5281/zenodo.23501588}. The
proof above is self-contained apart from the cited order-condition
framework; no private review is required to interpret it.

\textbf{Authorship and AI use.} The scholarly creator is Anonymous under
the Evidence Press publication convention. AI agents assisted with
mathematical development, literature inspection, drafting and internal
adversarial checks. Evidence Press is the publisher; account ownership
is not a claim of human research authorship. Internal AI review is not
external human peer review.

\textbf{Ethics, funding and competing interests.} This is theoretical
mathematical research without human participants or personal data. No
external research funding or competing interest was declared in the
supplied materials. These declarations describe the available record
rather than asserting an independently audited absence.

\section{Appendix A. Finite integer
certificates}\label{appendix-a.-finite-integer-certificates}

For each pattern, columns list the distinct words in lexicographic
order. Rows are the equations listed in Lemma 5, in that order. A row
\(P(a)\) has coefficient one in every word starting with a and zero
elsewhere; a row \(S(x,y)\) counts the interleavings of x and y,
retaining multiplicities.

For ab, use columns \((ab,ba)\): \[
\begin{pmatrix}1&0\\0&1\end{pmatrix},\qquad \det=1.
\] For aab, use columns \((aab,aba,baa)\): \[
\begin{pmatrix}
1&1&0\\0&0&1\\2&1&0
\end{pmatrix},\qquad \det=1.
\] For abc, use columns \((abc,acb,bac,bca,cab,cba)\): \[
\begin{pmatrix}
1&1&0&0&0&0\\
0&0&1&1&0&0\\
0&0&0&0&1&1\\
1&0&1&1&0&0\\
1&1&0&0&1&0\\
1&1&1&0&0&0
\end{pmatrix},\qquad \det=-1.
\] For aaab, use columns \((aaab,aaba,abaa,baaa)\): \[
\begin{pmatrix}
1&1&1&0\\0&0&0&1\\3&1&0&0\\3&2&1&0
\end{pmatrix},\qquad \det=1.
\] For aabb, use columns \((aabb,abab,abba,baab,baba,bbaa)\): \[
\begin{pmatrix}
1&1&1&0&0&0\\
0&0&0&1&1&1\\
2&1&1&0&0&0\\
2&1&0&1&0&0\\
0&1&0&2&1&0\\
4&2&0&0&0&0
\end{pmatrix},\qquad \det=2.
\] For aaaab, use columns \((aaaab,aaaba,aabaa,abaaa,baaaa)\): \[
\begin{pmatrix}
1&1&1&1&0\\0&0&0&0&1\\4&1&0&0&0\\6&3&1&0&0\\4&3&2&1&0
\end{pmatrix},\qquad \det=1.
\]

For the weight-seven boundary aaabb, the columns are \[
(aaabb,aabab,aabba,abaab,ababa,abbaa,baaab,baaba,babaa,bbaaa).
\] Every polarised-tail and proper-shuffle row annihilates the nonzero
column vector \[
(0,1,-1,-3,4,-1,2,-3,1,0)^T.
\] Conversely, take the nine rows \[
P(a),P(b),S(a,aabb),S(aa,abb),S(aaab,b),
S(a,abab),S(aa,bab),S(aab,ab),S(aaba,b).
\] Deleting the second column gives a nine-by-nine minor of determinant
one. These facts certify rank exactly nine. They concern the determining
equations, rather than the attainability of that remaining direction by
Williamson coefficients.

\section{References}\label{references}

Bazavov, A. (2022). Commutator-free Lie group methods with minimum
storage requirements and reuse of exponentials. \emph{BIT Numerical
Mathematics, 62}, 745--771.
\url{https://doi.org/10.1007/s10543-021-00892-x}

Bazavov, A. (2025). \emph{2N-storage Runge--Kutta methods: c-reflection
symmetry and factorization of the Butcher tableau} {[}Preprint{]}.
arXiv. \url{https://arxiv.org/abs/2506.08444}

Bazavov, A. (2026). 2N-storage Runge--Kutta methods: Order conditions,
general properties and some analytic solutions. \emph{BIT Numerical
Mathematics, 66}, Article 38.
\url{https://doi.org/10.1007/s10543-026-01139-3}

Lystad, T. A. (2026). \emph{All-stage order-three transfer for
Williamson 2N-storage Runge-Kutta methods} (Version 1.0.1)
{[}Preprint{]}. Zenodo. \url{https://doi.org/10.5281/zenodo.21764741} .
Archive metadata:
\url{https://github.com/arex1337/williamson-order-three/blob/main/CITATION.cff}
. Inspected manuscript source:
\url{https://github.com/arex1337/williamson-order-three/blob/main/paper/order-three-transfer.tex}

Owren, B. (2006). Order conditions for commutator-free Lie group
methods. \emph{Journal of Physics A: Mathematical and General, 39},
5585--5599. \url{https://doi.org/10.1088/0305-4470/39/19/S15} .
Inspected preprint:
\url{https://www.math.ntnu.no/preprint/numerics/2005/N7-2005.pdf}

Williamson, J. H. (1980). Low-storage Runge--Kutta schemes.
\emph{Journal of Computational Physics, 35}(1), 48--56.
\url{https://doi.org/10.1016/0021-9991(80)90033-9}

Yan, Y.-A. (2017). Low-storage Runge--Kutta method for simulating
time-dependent quantum dynamics. \emph{Chinese Journal of Chemical
Physics, 30}(3), 277--286.
\url{https://doi.org/10.1063/1674-0068/30/cjcp1703025}

\end{document}
